% 第 8 章　万有引力：苹果、月球与潮汐
\chapter{万有引力：苹果、月球与潮汐}

\step{反常识开场}

如果你住在海边，或者去海边旅行时留意过，会发现一件奇怪的事：海水每天涨落两次，大约每十二小时多一点涨一次潮。更奇怪的是，当月亮高高挂在头顶、把“正下方”的海水拉起来的时候，地球背对月亮那一侧的海水，居然也在涨潮。

请停下来想一想这件事有多反常。我们很容易接受“月球把正下方的海水吸起来”这个画面：引力嘛，隔着真空也能拉东西，拉最近的海水最用力，那里就鼓起来了。可是地球另一侧的海水，离月亮最远，被月球拉得最轻，它凭什么也鼓起来？如果引力真的只是“月球拽海水”，那最远端的海水应该被拽得最扁才对，那里应该是一天中水位最低的地方。但事实恰恰相反：最远端和最近端几乎同时涨潮，水位最低的地方反而是这两个方向的中间。

这个矛盾不是文字游戏，而是一条任何人都能核实的观测事实：打开任何一个港口城市的潮汐时刻表，你都会看到每天两次高潮，间隔大约十二小时二十五分钟，而不是一次。月球绕地球一圈要将近一个月，地球自转一圈是一天，怎么也算不出“两次”来——除非你承认，我们平时对引力的想象里缺了点什么。

这一章就是要补上这个“什么”。我们会发现，要解释背月一侧的涨潮，需要把引力从“月球拽东西”升级成一种更精确的想法：引力是一种随距离变化的拉扯，而真正掀起潮汐的，是这种拉扯在地球这一头和那一头的差别。而“差别”这个词，将成为本章后半段的主角。

顺便说一句，本章开场用的是潮汐，但最后我们解释的不只是潮汐。空间站里漂浮的宇航员、天上不偏不倚悬着的通信卫星、被甩出太阳系的旅行者号探测器，甚至你手里这部手机能定位到几米以内，背后全是同一套想法。

\step{先猜一猜}

在往下看之前，请先诚实地写下你对下面几个问题的直觉答案。不用查资料，猜错了才好玩。

第一题：月球正下方的海水涨潮，地球背对月球那一侧的海水也涨潮。你觉得原因是：（a）月球的引力绕过地球传到了对面；（b）地球被月球拉得比远端海水更厉害，相对地“离开”了远端海水；（c）月球把地球大气压过去，把海水挤起来了。

第二题：空间站里的宇航员会飘起来，你猜是因为：（a）那里太高，地球引力已经小到可以忽略；（b）那里是真空，没有空气托着所以飘着；（c）别的原因。

第三题：从一座高塔上水平扔出一块石头，扔得越用力，石头飞得越远才落地。假如你力气可以无限大，扔出的石头最终会：（a）直飞冲天，永远离开地球；（b）绕地球一圈飞回来砸中你的后脑勺；（c）无论多用力都会落到地上，只是落点更远。

第四题：你觉得“苹果落地”和“月球绕地球转”这两件事，背后是同一种规律，还是两种不同的规律？

先记住你的答案。这一章结束时，我们会回来逐条对答案。如果你第二题选了（a），你并不孤单——绝大多数人都是这么想的，而它恰好错得最典型。

\step{动手实验}

本章要建立的最重要的直觉是：一个向四面八方均匀散开的东西，越远越稀薄。引力就是这样散开的。这个直觉可以用一把手电筒和一部手机在家里亲自摸到。

\begin{experiment}[光斑为什么会变大变暗]
关掉房间的灯，打开手机的手电筒（或者任何一把手电），对着白墙照出一个光斑。让手机贴在墙上，然后慢慢向后退。

你会看到两件事同时发生：光斑越来越大，也越来越暗。如果手电能保持光斑大致是圆的，你可以粗略地量一量：后退到原来两倍远时，光斑的直径大约变成两倍，面积大约变成四倍，而光斑的亮度大约只有原来的四分之一。再退到三倍远，面积大约九倍，亮度大约九分之一。

如果你的手机能装一个测光强的应用（很多免费应用可以读取环境光传感器），把手机固定在一处，让手电筒从近到远移动，记下读数。你会发现亮度大体上按“距离平方分之一”往下掉：两倍远剩四分之一，三倍远剩九分之一。

为什么？因为手电发出的光总量没有变，只是摊在了越来越大的面积上。光从一个小灯泡出发，沿直线向四面八方散开，像吹气球时气球表面那样铺开：气球的半径变成两倍，整个球面的面积就变成四倍，原来涂在一块补丁上的颜料，现在要涂满四倍大的补丁，颜色自然淡了四倍。

这个实验像什么、不像什么，得说清楚。它像引力的地方在于：引力和光一样从一个源头向所有方向均匀散开，所以它们的“稀薄程度”都按距离平方分之一变化。它不像引力的地方在于：光是可以被墙挡住、被吸收的能量流，而引力几乎无法屏蔽——你没法在月球和海水之间插一块“引力挡板”。实验只借用“散开”这个几何事实，别借得太多。
\end{experiment}

还有第二个小体验，只需要一根一米左右的结实绳子和一块橡皮。把橡皮拴在绳端，在头顶上慢慢甩，让它做圆周运动，然后逐渐加快。你会清楚地感觉到：甩得越快，绳子对手的拉力越大；绳子放长一些、保持同样的转速，拉力也会变化。你的手在做的，就是地球对月球做的事情——通过绳子不断把橡皮“往中间拽”，橡皮才没有沿直线飞走。请把“转圈需要有人不停地拽”这个手感记住，它是本章理解轨道的钥匙。

\step{旧模型遇到麻烦}

现在回头看看，在没有万有引力这个概念之前，人们是怎样解释世界的，以及那种解释为什么走到头了。

最古老也最自然的想法是：天上的东西和地上的东西服从两套不同的规矩。地上的东西——石头、苹果、水——天然地往低处掉，掉到地面就停了，这是“重物的本性”；天上的东西——月亮、行星、星星——不会掉下来，它们天生就是转圈的，圆周运动是“天体的本性”。这套说法统治了人类将近两千年，因为它和日常观察贴得太紧了：你从没见过月亮掉下来，也没见过苹果绕着什么转圈。

但这个“两套规矩”的模型有一个致命的软肋：它没有回答“为什么”。为什么天体就该转圈？为什么地上物体就该下落？如果追问下去，答案只能是“因为它们本来就这样”。这在物理学里不算解释，只算把问题换了个说法。更麻烦的是，它无法做任何定量的预测：它说不出月亮转一圈该用多少天，也说不出潮水一天该涨几次。

第二个麻烦来自一个思想实验，你可以现在就做。站在椅子上（站稳，或者想象也行），一手拿一块小石子，水平扔出去。石子飞出几米，划一道弧线落地。再想象你力气大一倍，石子飞得远一倍。继续想：力气大到石子能飞出一公里、十公里、一百公里——这时会发生一件微妙的事：地球是圆的，石子在往前飞的同时，地面在它脚下往下弯。如果石子飞得足够快，它每往前飞一段，地面就恰好往下弯掉同样的高度，石子永远“接近地面”却永远“踩不到地面”。它会一直下落，却一直落不下来——它绕着地球转起圈来了。

这就是牛顿著名的“山顶大炮”思想实验：在高山之巅架一门大炮，水平打出的炮弹速度越快，落点越远；速度大到某个值，炮弹就再也落不下来了。注意这个画面里最关键的一句：炮弹不是“摆脱了下落”，而是“一直在下落”。它绕地球转的每一秒钟，都在向地面掉，只是地面在躲它。

\begin{center}
\begin{tikzpicture}[scale=1.05]
  \fill[green!12] (0,0) circle (1.4);
  \draw (0,0) circle (1.4);
  \fill (0,1.4) circle (1.2pt);
  \node[above=2pt] at (0,1.4) {山顶大炮};
  \draw[->,thick] (0,1.4) .. controls (0.85,1.35) .. (30:1.42);
  \draw[->,thick] (0,1.4) .. controls (1.35,1.42) and (1.7,0.85) .. (-12:1.45);
  \draw[dashed,thick] (0,0) circle (1.6);
  \node at (0.55,1.05) {慢};
  \node at (2.05,1.1) {更快};
  \node at (1.6,-1.35) {恰好永远落不下来};
\end{tikzpicture}
\end{center}

于是“两套规矩”模型被撞开了一条缝：只要速度够大，一颗炮弹——一个不折不扣的“地上物体”——也能像月亮一样转圈。反过来说，月亮也许就是一颗一直在下落的炮弹。天上和地下之间那堵墙，开始晃了。

还有第三个麻烦，来自精确的天文观测。早在牛顿之前，开普勒就整理了第谷几十年测火星的数据，发现行星轨道根本不是完美的圆，而是椭圆；行星走得时快时慢，离太阳近时快、远时慢；而且所有行星的周期和轨道大小之间有一种严格的比例关系。“天体天生做完美圆周运动”这个古老信条，其实早就被数据推翻了，只是当时没有人能回答：那到底是什么力量让行星走椭圆、还走得忽快忽慢？

三个麻烦指向同一个方向：我们需要一个新概念，把苹果和月亮、把下落和转圈、把天上和地下，装进同一套规律里。

\step{发明新概念}

牛顿的天才不在于他“发现了引力”——谁都知道东西会往下掉。他的天才在于提出了一个大胆得近乎冒失的断言：让苹果落地的那种拉扯，和让月亮绕地球转的那种拉扯，是同一种拉扯；而且不只地球有，月球有，太阳有，每一块石头都有。任何两个物体，无论是什么，隔着多远，都互相拉扯。这就是“万有引力”里“万有”两个字的全部含义：不是“地球上普遍有”，而是“宇宙里普遍有”。

这个断言立刻面临一个定量的问题：这种拉扯随距离怎么变？牛顿的答案是平方反比——距离变成两倍，拉扯弱成四分之一；距离变成十倍，弱成百分之一。为什么是平方，而不是一次方或者三次方？这里有一个干净的几何理由，正是你手电实验里看到的：想象从引力源出发的“拉扯的总量”向四面八方均匀散开，它要铺满一个又一个以源头为中心的球面。球面半径是 \(r\) 时，面积是 \(4\pi r^{2}\)，同样一份总量摊在四倍大的球面上，每一点分到的自然就弱成四分之一。平方反比不是凑出来的数字，它是“三维空间里均匀散开”的必然结果。我们生活的空间恰好是三维的，所以引力恰好是平方反比——如果空间是四维的，引力就会按三次方衰减。

\begin{center}
\begin{tikzpicture}[scale=1.0]
  \fill (0,0) circle (1.5pt);
  \node[left] at (0,0) {源};
  \draw (0,0) -- (24:4.6);
  \draw (0,0) -- (-24:4.6);
  \draw (24:2) arc[start angle=24,end angle=-24,radius=2];
  \draw (24:4) arc[start angle=24,end angle=-24,radius=4];
  \node at (2.6,0) {\(S\)};
  \node at (4.85,0) {\(4S\)};
  \node[below] at (1.95,-1.05) {距离 \(r\)};
  \node[below] at (3.9,-1.35) {距离 \(2r\)};
\end{tikzpicture}
\end{center}

同样一份“总量”，在距离 \(r\) 处摊在大小为 \(S\) 的接收面上；距离加倍，同样的射线散开成横竖各两倍的面积 \(4S\)，每单位面积分到的只剩四分之一。

这个比喻像什么、不像什么，也要交代清楚。说引力像“摊开的光”，像的是衰减的几何；但光是有东西真的在飞（光子），而牛顿引力里并没有任何“引力子弹”在飞行，它更像空间中每一点都预先写好了一张“此地该受多大拉扯、朝哪个方向”的告示牌。这张遍布空间的告示牌，就是我们发明的新概念：引力场。物体不是因为“隔空认出”了另一个物体才受力，而是因为它读到了自己所在位置的那张告示牌。场这个概念此刻也许显得多余，但请记住它——到了电磁学那几章，它会从“方便的记账方式”变成“必须存在的实体”，而在更后面的章节里，连引力场本身也会被改写。

有了“互相拉扯、平方反比”，还需要最后一块拼图：拉扯产生的后果是什么。答案在前几章已经备好了——力改变的是动量，是运动状态；不是“物体需要力才能维持运动”，恰恰相反，没有力物体就沿直线匀速飞。于是月亮为什么不掉下来这个问题，被彻底翻转了：月亮其实一直在掉。它本来有沿直线飞走的趋势，地球的引力每时每刻把它往地心方向拽一点，拽出的轨迹恰好弯成一个闭合的圈。转圈不是月亮的“本性”，是月亮在引力和惯性之间达成的动态平衡。同理，你手里甩着的橡皮也不是“想转圈”，是绳子在不停地拽它。这里要特别说清楚：所谓“向心力”不是一种和重力、弹力并列的新种类的力，它只是一个角色——谁的合力恰好指向圆心、让物体转弯，谁就扮演这个角色。对月亮，这个角色由引力出演；对橡皮，由绳子的拉力出演。

\step{一句话模型}

\begin{oneline}
任何两个有质量的物体都互相吸引，力的大小与质量的乘积成正比、与距离的平方成反比；天体不是在“绕圈”，而是在不停地自由落体，只是一直错过地面——而潮汐，是这种吸引在物体两端强弱不同撕出来的。
\end{oneline}

\step{数学翻译器}

现在把上面的每一句话翻译成可以计算的式子。

第一句话——“互相吸引、平方反比”——写成
\[
F = G\,\frac{m_{1} m_{2}}{r^{2}},
\]
其中 \(m_{1}\)、\(m_{2}\) 是两个物体的质量，\(r\) 是它们中心之间的距离，\(G\) 是一个由实验测定的小常数，\(G \approx 6.67\times 10^{-11}\,\mathrm{N\cdot m^{2}/kg^{2}}\)。先看这个式子回答什么：左边是“两个物体各自受到多大的拉扯”；右边每一项都在说“什么让拉扯变强或变弱”——质量越大拉扯越强（正比），距离越远拉扯越弱（平方反比），\(G\) 负责把单位和宇宙的实际强度对上。

拿到公式立刻做特殊情形检查。距离变成两倍：分母变成四倍，力弱成四分之一——和手电实验一致，通过。距离趋于无穷远：力趋于零——两个离得无限远的物体互不影响，合理，通过。质量变成零：力变成零——没有质量就没有参与拉扯的资格，合理。两个物体的位置对调：力大小不变、方向相反——这正是作用与反作用，公式自动满足，通过。还有一个值得多想的检查：\(r\) 趋于零时力趋于无穷大，这看着吓人，但其实暴露了公式的一个隐藏条件——它把物体当成了“质点”。真实物体有大小，两个篮球没法让中心距离等于零，所以这个“无穷大”在现实中永远碰不到；但引力源真的能被压缩到很小吗？先留个悬念，到“模型的边界”那一步再拆。

第二个翻译是把“月亮一直在下落”变成数字，这是物理史上最漂亮的一次核对。地面上，苹果下落的加速度是 \(g \approx 9.8\,\mathrm{m/s^{2}}\)。地心到苹果的距离约等于地球半径 \(R \approx 6.4\times10^{6}\,\mathrm{m}\)，而地心到月球的距离大约是 \(60 R\)。如果引力真是平方反比，月球位置上由地球引力产生的加速度应该是
\[
a_{\text{月}} = \frac{g}{60^{2}} = \frac{9.8}{3600} \approx 2.7\times10^{-3}\,\mathrm{m/s^{2}}.
\]
另一边，月球绕地球一圈约 \(27.3\) 天，轨道半径 \(60R\)，由圆周运动的向心加速度公式 \(a = 4\pi^{2}r/T^{2}\) 算出来的加速度，也是 \(2.7\times10^{-3}\,\mathrm{m/s^{2}}\) 左右。两个完全独立的算法——一个从“苹果落地”外推，一个从“月亮的圈有多大、转多快”实测——给出了同一个数。天上的运动和地上的运动，第一次被一个数字缝在了一起。据说牛顿当年算出这个结果后，把文稿压了将近二十年才发表，因为他还缺一块拼图：地球这么大，凭什么能把它的全部质量当成集中在地心来算？

\begin{mathbridge}
这块拼图叫壳层定理：一个质量均匀球对称分布的球体，对球外任何一点的引力，精确等于把全部质量集中在球心时产生的引力；而球壳内部的点，感受不到球壳的引力。证明需要对球面上每一小块的引力做积分：离得近的小块拉得狠但数量少，离得远的小块数量多但拉得轻，积分之后所有方向的分量相互抵消或叠加，最后恰好拼出“等效于质点”这个结果。平方反比在这里是关键——只有平方反比（以及它的“镜像”，正比于距离的胡克力）能让这个奇迹发生。这也是为什么我们可以心安理得地用地球半径去算月球、用地心距离去算苹果。

用同样的平方反比，还能推出圆轨道上的速度。设卫星绕地球做半径为 \(r\) 的圆周运动，引力提供向心作用：
\[
G\frac{M m}{r^{2}} = m\frac{v^{2}}{r}
\quad\Longrightarrow\quad
v = \sqrt{\frac{GM}{r}}.
\]
注意卫星自己的质量 \(m\) 被约掉了：同一高度上，无论卫星是一公斤还是一吨，绕行速度都一样。这其实是“所有物体下落一样快”的轨道版。近地轨道（\(r\) 比地球半径略大）代入数字，\(v \approx 7.9\,\mathrm{km/s}\)——这就是“第一宇宙速度”，那门山顶大炮需要的炮口速度。每秒八公里，差不多是步枪子弹速度的十倍，这就是为什么人类到二十世纪才把卫星送上天。
\end{mathbridge}

第三个翻译给开普勒的三条定律一个家。开普勒从数据里归纳出：（一）行星轨道是椭圆，太阳在椭圆的一个焦点上；（二）行星和太阳的连线在相等时间内扫过相等的面积——所以行星离太阳近时跑得快、远时跑得慢；（三）各行星周期的平方与轨道半长轴的三次方成正比，\(T^{2} \propto a^{3}\)。牛顿证明了这三条全是平方反比引力的推论：椭圆轨道来自平方反比力场的精确求解；面积定律来自引力的方向永远指向中心、因而不会改变行星对中心的角动量；第三条则可以由圆轨道的近似一眼看出来——把上面的 \(v = \sqrt{GM/r}\) 代进 \(T = 2\pi r/v\)，平方之后就得到 \(T^{2} = 4\pi^{2}r^{3}/(GM)\)。数据归纳和经验公式，第一次被还原成了更基本规律的模样。

这条 \(T^{2} \propto r^{3}\) 关系有一项每天都在用的工程应用：通信和导航卫星。想让一颗卫星永远悬在赤道上空某一点不动（同步卫星），它的周期必须等于地球自转周期 \(T = 24\) 小时。把近地卫星的数据（周期约九十分钟、轨道半径约一个地球半径）按比例外推：周期变成 \(24\) 小时是九十分钟的 \(16\) 倍，周期的平方是 \(256\) 倍，所以轨道半径的三次方也要 \(256\) 倍，半径约为 \(256^{1/3} \approx 6.3\) 倍——即约六点三个地球半径、四万公里上下的轨道半径，减去地球半径，就是三万六千公里上下的高度。这正是真实同步卫星的轨道高度，电视锅盖天线永远对准同一个方向，靠的就是它。反过来，把这套比例用到行星际航行，还能解释“引力弹弓”：探测器飞掠行星时借行星公转的动量加速，省下大量燃料——旅行者号正是这样一站接一站地被甩出太阳系外围的。引力理论在这里不是考试题，而是航天工程师的施工图。

第四个翻译，终于轮到开场的主角：潮汐。月球对地球的引力随距离平方反比变化，这意味着它对地球“近月一侧”的海水、对地球中心、对“远月一侧”的海水，拉的力度是三个不同的数。近端拉得最狠，地心次之，远端最轻。以地心为参照来看：近端海水被拉得比地心多，于是它向月球方向隆起；远端海水被拉得比地心少，于是它“落后”于地心，向远离月球的方向隆起。两边一起涨。这种“同一物体两端受力不同”的差值，叫作潮汐力。注意它的大小：引力的差别随距离衰减得比引力本身更快——引力按 \(1/r^{2}\)，潮汐力按 \(1/r^{3}\)。这解释了为什么质量只有太阳 \(1/27000000\) 的月球，对地球潮汐的贡献反而是太阳的两倍多：月球近得多，而潮汐力对距离的敏感多出一个平方。

\begin{challenge}
潮汐力的定量形式可以这样推：设月球质量 \(M\)，地月中心距离 \(d\)，地球半径 \(R\)。月球引力在地心处产生的加速度为 \(GM/d^{2}\)，在近月点为 \(GM/(d-R)^{2}\)。把后者对 \(R\) 展开（\(R \ll d\)，只保留一阶项），两者之差约为 \(2GMR/d^{3}\)。远月点算出来大小相同、方向相反。因子 \(R/d^{3}\) 就是“潮汐力正比于被撕物体尺寸、反比于距离三次方”的来源。这个 \(1/d^{3}\) 的依赖性在极端天体附近会变得极其重要：它说明越靠近引力源、或者自身尺寸越大，被“意面化”撕裂的风险越高。还能顺手解释一个观测：地球对月球的潮汐作用早就把月球的自转“锁定”成和公转同步，所以月球永远以同一面朝向地球——这不是巧合，是潮汐摩擦几十亿年磨出来的结果。
\end{challenge}

\step{模型的边界}

一个好的模型必须能说清自己在哪里会错，牛顿引力在这方面堪称模范。

第一，它假设物体是球对称的（或可当质点处理）。对行星和卫星这近似极好，但对“地球表面的重力随纬度变化”这类精细问题就不够了：地球自转让它赤道略鼓，赤道上的引力既被略大的半径削弱，又被自转的惯性效应分走一小份，所以赤道的 \(g\) 比两极小约百分之零点五。这不是牛顿定律的失败，而是“均匀球”这个近似该退休的地方。

第二，它是“平方反比、瞬时、超距”的。牛顿的公式里，引力传递不需要时间：太阳此刻消失，地球立刻沿直线飞出去。牛顿本人对此深感不安，他在信里写道，没有介质却能隔空传力“荒谬到任何有哲学头脑的人都不能接受”，但他选择“我不杜撰假设”——先让公式干活，机制留给后人。这个边界在两百年后被爱因斯坦补上：引力不是瞬时的拉力，它以光速传播；它不是“力”，而是时空弯曲的表现。在日常尺度上，两种理论算出的差别小得几乎测不出——除了一个著名的例外：水星。水星离太阳最近、引力最强，它的轨道近日点每世纪多出 \(43\) 角秒的进动，牛顿力学把其他行星的拉扯全算上还差这 \(43\) 角秒，广义相对论精确补上了它。这告诉我们牛顿引力的适用范围：弱引力场、远小于光速的运动。靠近黑洞、研究宇宙大尺度结构时，必须换理论。

第三，最典型的误用清单。“太空里没有引力，所以宇航员失重”——错。空间站高度约四百公里，那里的引力强度仍有地面的九成左右；宇航员失重是因为空间站和宇航员一起在做自由落体，一起“往下掉”，谁也不需要支撑谁。“重力加速度就是 \(9.8\)”——那只是地表附近的近似值，高度、纬度、地下矿藏都会让它变化，地质勘探正是靠测这种微小变化来找矿的。“引力平方反比，所以靠近引力源就能无穷大”——前面说过，质点近似在物体内部失效，壳层定理保证了钻进地球内部时引力反而随深度线性减弱。还有“月球对潮汐的作用是把它正下方的海水吸起来”——读完本章你知道，这是把“平均的拉”误当成了“差别的拉”。

\step{回头解释开场现象}

现在回到海边。地球和地球上的海水组成一个被月球（还有太阳）拉扯的系统。月球引力在空间中的分布是：近月点最强，地心次之，远月点最弱。海水是流体，可以响应这个差别；固体地球整体近似以地心处的力度被拽着走。于是在随地球中心一起运动的参考系里看：近月一侧的海水“被多拽了一把”，隆起指向月球；远月一侧的海水“被少拽了一把”，被地球主体甩在后面，隆起背向月球。两个隆起，两个低潮区，地球每天在这副固定的隆起图案下面自转一周，所以每个沿海地点一天穿过两次高潮——开场时刻表上的“十二小时二十五分钟”，就是地球自转周期加上月球绕行的微小修正。

\begin{center}
\begin{tikzpicture}[scale=1.0]
  \fill[blue!15] (0,0) ellipse (2.1 and 1.35);
  \fill[green!20] (0,0) circle (1.05);
  \draw (0,0) circle (1.05);
  \fill[gray!50] (4.6,0) circle (0.42);
  \node at (4.6,-0.85) {月球};
  \draw[->,very thick] (1.5,0) -- (2.9,0);
  \node[above] at (2.2,0.02) {近端拉得最狠};
  \draw[->,thick] (0.1,0) -- (1.05,0);
  \draw[->,thick] (-1.5,0) -- (-0.75,0);
  \node[above] at (-1.6,0.15) {远端拉得最轻};
  \node[below] at (0,-1.5) {两端海水相对地心都隆起};
\end{tikzpicture}
\end{center}

图中三支箭头是月球引力在近月点、地心、远月点三处的强弱：都指向月球，但一支比一支短。流体响应的是这个“长短之差”，而不是箭头本身——所以隆起出现在两端，而不是只出现在月球正下方。

顺便回答那几个猜一猜。第一题答案是（b）：不是引力“绕过去”，而是地球被拉得比远端海水更狠。第二题，如果你选了（a），现在可以改过来了：空间站处引力并不小，失重的真正原因是自由落体——宇航员、空间站、空间站里的水和笔，全都在以同样的加速度一起下落，于是彼此之间呈现“漂浮”。这也正是为什么抛物线飞行的飞机能在机舱里制造几十秒的失重训练环境，以及为什么跳楼机下坠的瞬间你会觉得胃往上提。第三题：力气无限大时石头绕地球飞回来——理想情况下它真的可能砸中你的后脑勺（忽略空气阻力），更可能的情况是速度足够大时它再也回不来，沿着一条张开的曲线离开地球；那个临界速度，由引力决定，叫作逃逸速度，地面上约每秒十一公里。第四题：同一种规律，这正是全章的核心。

还有一个日常的回音值得提一句：为什么潮汐里太阳也有份？满月和新月时，太阳、月球几乎在一条线上，两者的潮汐隆起叠加，潮差最大，叫大潮；上弦月和下弦月时两者方向垂直，隆起互相抵消一部分，潮差最小，叫小潮。渔民和赶海人的潮汐表里那套“初一十五大潮”，就是这么来的。一套公式同时解释了苹果、月亮、炮弹、卫星和渔汛，这就是“万有”二字的分量。

\step{脑洞实验与挑战题}

本章的最后一个想法，是通向下一座山的桥。爱因斯坦后来回忆说，他一生最快乐的一个念头是：一个从屋顶上自由落下的人，在下落过程中感觉不到自己的重量。把这句话展开成一个思想实验：

\begin{thought}[看不见窗外的电梯]
想象你在一部密闭电梯里，脚下踩着一台体重秤。电梯静止在地面时，秤显示你的体重。现在缆绳剪断，电梯和你一起在地球引力场中自由下落——秤的读数立刻变成零，你飘了起来，和空间站里的宇航员一模一样。反过来，如果这部电梯被放到没有任何引力的深空，用火箭以 \(9.8\,\mathrm{m/s^{2}}\) 的加速度向上推，你脚下的秤又会显示出你平时的体重，你手里的苹果松手后照样“落”向地板。

关键的问题是：不看窗外，你做任何物理实验，能区分这两种情况吗——“引力中的自由落体”和“无引力的漂浮”，“静止在地面”和“被火箭加速”？牛顿力学里找不到任何实验能区分它们。爱因斯坦把这个观察上升为一条原理——等效原理：在足够小的时空范围内，引力的效果和加速运动的效果不可区分。沿着这条路走下去，引力将不再是“力”，而是时空本身的几何。那是后话的章节。这里先记住它的入口，正是本章反复出现的事实：所有物体无论质量，在引力中下落得一样快；以及失重不是“没有引力”，而是“跟着引力一起掉”。
\end{thought}

这套想法并不只停在哲学层面，它每天都在你的手机里干活。导航卫星在约两万公里的高空，那里的引力比地面弱、卫星又以每秒近四公里的速度飞行，按狭义和广义相对论，卫星上的原子钟每天会比地面快约三十八微秒。听起来微不足道？光是每秒三十万公里，三十八微秒对应十多公里的定位误差——如果不把本章的引力理论和相对论的修正一起写进卫星的时钟程序，你的导航每天会漂出十几公里。一颗卫星就这样把引力、轨道、原子钟和相对论串在了一起，这也是后文“一颗卫星”这位老朋友要反复登场的预告。

\begin{puzzles}
\item 假设挖一条笔直穿过地心的隧道，忽略空气阻力和地热，你在隧道口松手跳进去，接下来会怎样？请定性描述整个运动过程。（思路提示：壳层定理说，当你身处地球内部、离地心距离为 \(r\) 时，只有你脚下那个半径为 \(r\) 的球的质量在拉你，外层球壳的贡献全部抵消。地球密度近似均匀时，这部分质量正比于 \(r^{3}\)，再除以 \(r^{2}\)，拉力正比于 \(r\)。你见过“拉力正比于位移、方向指向中心”的运动。）

\item 天文观测发现，月球正以每年约 \(3.8\) 厘米的速度慢慢远离地球，同时地球的自转在变慢，一天在变长。这里的“变慢”和“远离”是同一件事的两面，请用潮汐解释：地球自转的什么东西转移给了月球？（思路提示：海水隆起被地球自转拖着跑，并不严格指向月球，而是略微“超前”；这团偏了的隆起对月球的拉力也跟着偏了，不再恰好指向地心。想想这个偏了的拉力对月球起什么作用，以及它对地球的反作用。守恒的总量是地月系统的角动量。）

\item 地球同步卫星悬在赤道上空约三万六千公里处，与地面保持相对静止。为什么它必须在赤道上空，而不能悬在北京上空？如果地球自转突然变慢一倍，同步轨道的高度会怎么变？（思路提示：卫星做圆周运动所需的指向地心的作用只能由引力出演，而引力永远指向地心——轨道圆心必须是地心。高度变化用 \(T^{2}\propto r^{3}\) 估，不用算出精确值，说出变大还是变小、大约变到几倍。）

\item 靠近黑洞的宇航员会被潮汐力沿头脚方向拉长，科幻里叫“面条化”。有趣的是：掉进超大质量黑洞（比如星系中心那种）时，宇航员在跨过“视界”时可能几乎毫无感觉；而掉进一个只有几倍太阳质量的小黑洞，还没到视界就会被撕碎。质量更大的黑洞，为什么反而更“温柔”？（思路提示：潮汐力正比于 \(M/d^{3}\)，而黑洞视界的半径正比于 \(M\)。把视界处的潮汐力用 \(M\) 表示出来，看看它随 \(M\) 增大是变大还是变小。）
\end{puzzles}

